\section{Exhaustive enumeration and finite determination of transport operators}
\label{iii:nucleo-finito}

This appendix develops two finite operations with different starting
data. The first enumerates the initial conditions of the two cursors
and constructs their emissions through the declared arithmetic recurrence.
The second determines transport operators from the regional
transition register. Separating their domains specifies
what each proof produces: enumeration constructs the emission
catalogue; matrix reconstruction determines the operators compatible
with the register. The reconstruction theorem is not used as a proof
of the prior production of that register.

The primitives of the first operation are the support $I_9^2$, the four
cardinal orientations, the tritic selector, the pre-displacement readout,
orientation reversals and the emission rule.
Their definitions are brought together below. Those of the second operation are the
four transition matrices in~\eqref{iii:nf-registro}.
Their unique reconstruction is proved in
Theorem~\ref{iii:nf-unicidad}; the resulting calendar is given
in~\eqref{iii:nf-calendario}. These determinations preserve
the starting data specified for each operation.

\subsection{Complete domain of initial conditions}

Let $I_9=\{0,\ldots,8\}$ and let
\[
 \mathcal D=\{N,E,S,O\},\qquad
 v_N=(-1,0),\quad v_E=(0,1),\quad
 v_S=(1,0),\quad v_O=(0,-1).
\]
An initial cursor condition is $s=(i,j,d)\in
\mathcal S=I_9^2\times\mathcal D$. The two cursors are independent
in the initial domain
\begin{equation}
 \Omega=\mathcal S_+\times\mathcal S_\times,
 \qquad |\mathcal S|=9^2\cdot4=324,
 \qquad |\Omega|=324^2=104\,976.
 \label{iii:nf-dominio}
\end{equation}
Index $i$ corresponds to the positive APP label $i+1$. For
positive integers, positive reduction is $\rho_9(n)=1+[(n-1)]_9$,
where $[a]_b$ denotes the representative in $\{0,\ldots,b-1\}$.
The residual evaluations of the two sheets are
\[
 \sigma(i,j)=\rho_9(i+j+2),\qquad
 \varpi(i,j)=\rho_9((i+1)(j+1)).
\]
The finite multiplicative readout uses
\begin{equation}
 \begin{array}{c|rrrrrrrrr}
 a&1&2&3&4&5&6&7&8&9\\ \hline
 \ell_9(a)&0&1&0&2&5&0&4&3&0
 \end{array}.
 \label{iii:nf-log}
\end{equation}
On units, $\ell_9$ is the discrete exponent in base $2$.
On $3,6,9$, the value $0$ constitutes the declared extension of the
emission observable; membership in that sector remains in the
trajectory register.

For $1\le t\le54$, define
\[
 r_t=1+[(t-1)]_9,
 \qquad
 \tau_t=
 \begin{cases}
 +1,&r_t\in\{1,4,7\},\\
 -1,&r_t\in\{2,5,8\},\\
 0,&r_t\in\{3,6,9\}.
 \end{cases}
\]
Each cursor preserves its position $z_t$ immediately before step
$t$ and its orientation vector $v_t$. The additive cursor is active
when $\tau_t=+1$; the multiplicative cursor, when $\tau_t=-1$. For each
of them, with activity indicator $a_t$, the update is
\begin{equation}
 z_{t+1}=[z_t+a_tv_t]_9,
 \qquad
 v_{t+1}=
 \begin{cases}
 -v_t,&t\in\{27,54\},\\
 v_t,&t\notin\{27,54\}.
 \end{cases}
 \label{iii:nf-actualizacion}
\end{equation}
The readout at step $t$ is taken at $z_t$, before displacement.
The neutral event preserves both positions. This convention fixes
unambiguously the window endpoints and the effect of
cardinal reversals.

In the windows $W_m=\{9m-8,\ldots,9m\}$, $1\le m\le6$, let
\begin{align}
 S_m(s_+)&=\sum_{\substack{t\in W_m\\\tau_t=+1}}
                 \sigma(z_t^+),&
 E_m(s_\times)&=\sum_{\substack{t\in W_m\\\tau_t=-1}}
                 \ell_9(\varpi(z_t^\times)),
 \label{iii:nf-firmas}\\
 s_m&=[S_m]_{10},&e_m&=[E_m]_{10}.\notag
\end{align}
Each window contains three events of each type. The signatures $s$ and $e$
are the six components of these residues, not the integers $S_m,E_m$
before reduction.

\subsection{Census of signatures and emissions}

The emission of a pair of initial conditions is determined by
\begin{equation}
 U_m=100s_m+10[s_m+e_m]_{10}+[3s_m+5e_m+1]_{10},
 \qquad 1\le m\le6.
 \label{iii:nf-emision}
\end{equation}
The final term $1$ is $[7\cdot3]_{10}$, corresponding to the three
neutral events. This rule's coefficients belong to the declared
recurrence. The real constants studied subsequently
do not enter enumeration of its domain or its evaluation.

\begin{proposition}[Injective factorization of the finite catalogue]
\label{iii:nf-inyectividad}
Emission distinguishes every signature pair $(s,e)$. There are $18$ additive
and $26$ multiplicative signatures; their pairs produce $468$ emissions.
The multiplicities of these emissions are $144$, $432$ and $1944$,
with $432$, $18$ and $18$ emissions, respectively.
\end{proposition}

\begin{proof}
The hundreds digit of $U_m$ recovers $s_m$. If $b_m$ is its
tens digit, then $e_m=[b_m-s_m]_{10}$. The units digit verifies
the third congruence of \eqref{iii:nf-emision}. Thus, the
coupling is injective.

The census is obtained by traversing, in the full Cartesian order,
$i=0,\ldots,8$, $j=0,\ldots,8$ and $d=N,E,S,O$, and applying the $54$
steps of \eqref{iii:nf-actualizacion}. The $324$ additive conditions
are distributed over the following $18$ signatures, each with $18$
preimages. In the tables, a sequence of six digits represents
the six ordered components of a signature.
\begin{center}
\begin{tabular}{ccc}
$122564$&$122898$&$212645$\\
$212988$&$221456$&$221889$\\
$456221$&$465898$&$546988$\\
$564122$&$645212$&$654889$\\
$889221$&$889654$&$898122$\\
$898465$&$988212$&$988546$
\end{tabular}
\end{center}
For the multiplicative cursor, the signatures and their multiplicities are
\begin{center}
\begin{tabular}{rr@{\qquad}rr}
Signature&Multiplicity&Signature&Multiplicity\\ \hline
$000000$&$108$&$555555$&$24$\\
$177393$&$8$&$555717$&$8$\\
$177555$&$8$&$555771$&$8$\\
$177771$&$8$&$555933$&$8$\\
$339555$&$8$&$717555$&$8$\\
$339771$&$8$&$717717$&$8$\\
$339933$&$8$&$717933$&$8$\\
$393177$&$8$&$771177$&$8$\\
$393393$&$8$&$771339$&$8$\\
$393555$&$8$&$771555$&$8$\\
$555177$&$8$&$933339$&$8$\\
$555339$&$8$&$933555$&$8$\\
$555393$&$8$&$933717$&$8$
\end{tabular}
\end{center}
These tables result from the finite sums \eqref{iii:nf-firmas}, with
domain, readout order and reversals already specified. Their sums
of multiplicities are $18\cdot18=324$ and
$24\cdot8+24+108=324$. The Cartesian domain
\eqref{iii:nf-dominio} permits every signature pair to be realized.
Injectivity gives $18\cdot26=468$ emissions and the
multiplicities
\[
 (18\cdot24,\ 18\cdot8)=(432,144),\qquad
 (18,\ 18\cdot24)=(18,432),\qquad
 (18,\ 18\cdot108)=(18,1944),
\]
where each pair indicates the number of emissions and multiplicity of their fibres.
Finally,
\[
 432\cdot144+18\cdot432+18\cdot1944=104\,976.
\]
Enumeration thus exhausts the complete domain of initial conditions.
\end{proof}

The oriented ternary reduction is
\[
 \Psi(U)=([-U_1]_3,\ldots,[-U_6]_3).
\]
Applied to the $18\cdot26$ emissions of the preceding tables,
it produces the following fibre census. Here $r$ counts distinct
emissions, not initial conditions.
\begin{equation}
 \begin{array}{c|rrrrrr}
 r&1&2&3&4&5&6\\ \hline
 \#\{w:\#\Psi^{-1}(w)=r\}&132&66&18&3&6&18
 \end{array}.
 \label{iii:nf-censo-ternario}
\end{equation}
The sum of the second row is $243$, while its weighted sum
with weights $r$ is $468$. The image of $243$ words is therefore distinguished
from the ambient space $\mathbb F_3^6$, containing $729$ words. The accompanying
certificate preserves all signatures with their preimages, the $468$
emissions and their ternary reduction; the preceding formulas provide
a complete procedure for reconstructing them.

\subsection{Transition register and reconstruction of the lifts}

This subsection fixes the regional register of the cited source.
Let $b_0^\chi,\ldots,b_4^\chi\in\mathbb F_3^6$, with
$\chi\in\{\mathsf C,\mathsf P,\mathsf A\}$, be its three sequences.
The register organizes transitions $b_0\to b_1\to b_2$ in
$X_0,Y_0$ and transitions $b_2\to b_3\to b_4$ in $X_1,Y_1$,
using row vectors and the same regional order:
\begingroup\small
\begin{equation}
\begin{aligned}
X_0&=\begin{pmatrix}
0&1&0&2&1&1\\0&1&2&2&2&2\\2&0&1&1&0&1\\
1&2&1&2&2&1\\1&2&1&2&0&0\\1&1&2&2&0&2
\end{pmatrix},&
Y_0&=\begin{pmatrix}
0&1&2&2&2&2\\0&1&0&2&1&1\\1&2&1&2&2&1\\
1&0&2&0&1&1\\1&1&2&2&0&2\\1&2&1&0&2&0
\end{pmatrix},\\[1ex]
X_1&=\begin{pmatrix}
0&1&0&2&1&1\\0&0&2&1&1&1\\1&0&2&0&1&1\\
0&1&2&2&2&2\\1&2&1&0&2&0\\0&1&0&2&1&0
\end{pmatrix},&
Y_1&=\begin{pmatrix}
0&0&2&1&1&1\\1&1&0&2&2&1\\0&1&2&2&2&2\\
1&0&2&0&1&1\\0&1&0&2&1&0\\0&1&0&2&0&0
\end{pmatrix}.
\end{aligned}
\label{iii:nf-registro}
\end{equation}
\endgroup
The preceding production is expressed in its source through
the composition of selection, transport and update:
\[
 x_{j+1}^{\chi}
 =\mathcal U_{9j+9}\circ\cdots\circ\mathcal U_{9j+1}(x_j^{\chi}),
 \qquad b_j^\chi=\Pi_6(x_j^\chi),
 \qquad\mathcal U_t=\mathsf{Upd}_t\circ\mathsf{Tra}_t\circ\mathsf{Sel}_t.
\]
The following result starts from register \eqref{iii:nf-registro}.
Its proof reconstructs the operators and verifies
transition compatibility; it does not attribute coordinatewise
production of states $x_j^\chi$ to matrix inversion.

\begin{proposition}[Uniqueness on the regional register]
\label{iii:nf-unicidad}
The systems $X_aL_a=Y_a$, $a=0,1$, have a unique solution over
$\mathbb F_3$. The two solutions are
\begingroup\small
\begin{equation}
L_0=\begin{pmatrix}
2&2&2&1&2&1\\2&1&2&2&1&1\\1&0&0&1&0&2\\
0&0&1&1&1&0\\2&2&2&0&2&1\\2&1&2&1&0&0
\end{pmatrix},\qquad
L_1=\begin{pmatrix}
2&1&2&0&2&2\\1&2&1&1&2&0\\1&2&1&1&1&2\\
0&1&0&0&2&1\\2&0&2&1&0&1\\0&2&2&2&1&1
\end{pmatrix}.
\label{iii:nf-lifts}
\end{equation}
\endgroup
\end{proposition}

\begin{proof}
Elimination in $\mathbb F_3$ gives $\det X_0=1$ and
$\det X_1=2=-1$. Each map $L\mapsto X_aL$ is therefore
an automorphism of the $36$-dimensional matrix space; under columnwise
vectorization it is represented by $I_6\otimes X_a$.
One obtains $L_a=X_a^{-1}Y_a$. Multiplication of the matrices in
\eqref{iii:nf-lifts} by the respective matrices $X_a$ reproduces the
two matrices $Y_a$ of \eqref{iii:nf-registro}; existence and uniqueness
are established on that register.
\end{proof}

The three resulting sequences, explicitly separated into blocks,
are
\begin{align*}
 B^{\mathsf C}&=010211\mid012222\mid010211\mid002111\mid110221,\\
 B^{\mathsf P}&=201101\mid121221\mid102011\mid012222\mid102011,\\
 B^{\mathsf A}&=121200\mid112202\mid121020\mid010210\mid010200.
\end{align*}
Each block preserves its six positions. Concatenation is not
identified with a base-ten integer.

\subsection{Exhaustiveness of the calendar and stability of reconstruction}

A binary calendar is $c=(c_0,c_1,c_2,c_3)\in\{0,1\}^4$.
Define the number of transitions incompatible with the register by
\begin{equation}
 d(c)=\sum_{\chi\in\{\mathsf C,\mathsf P,\mathsf A\}}
       \sum_{j=0}^{3}
       \mathbf1\{b_j^\chi L_{c_j}\ne b_{j+1}^\chi\}.
 \label{iii:nf-calendario}
\end{equation}
Multiplying the preceding sequences by
\eqref{iii:nf-lifts} gives
\[
 d(c)=3\,d_{\rm H}(c,(0,0,1,1)),
\]
where $d_{\rm H}$ is Hamming distance: at each of the
four positions, the alternative operator fails to satisfy all three
regional transitions. The complete census is
\begin{center}
\begin{tabular}{cr@{\qquad}cr@{\qquad}cr@{\qquad}cr}
$c$&$d(c)$&$c$&$d(c)$&$c$&$d(c)$&$c$&$d(c)$\\ \hline
$0000$&$6$&$0100$&$9$&$1000$&$9$&$1100$&$12$\\
$0001$&$3$&$0101$&$6$&$1001$&$6$&$1101$&$9$\\
$0010$&$3$&$0110$&$6$&$1010$&$6$&$1110$&$9$\\
$0011$&$0$&$0111$&$3$&$1011$&$3$&$1111$&$6$
\end{tabular}
\end{center}
Consequently, $0011$ is the unique calendar reproducing the
three complete sequences.

Sensitivity to matrix entries likewise admits an
exhaustive proof without approximations. A single-entry modification
replaces $L_a$ by $L_a+\lambda E_{ij}$, with
$a\in\{0,1\}$, $1\le i,j\le6$ and
$\lambda\in\mathbb F_3\setminus\{0\}$. There are
$2\cdot6\cdot6\cdot2=144$ such modifications. Since
$X_a$ is invertible,
\[
 X_a(L_a+\lambda E_{ij})-Y_a
 =\lambda X_aE_{ij}\ne0.
\]
Every modification destroys at least one transition of the
corresponding register. The executable certificate also enumerates
the $144$ modifications and preserves, for each, the number of
violated transitions.

The proved property is rigidity of the operators with respect to the
fixed regional register. Reproducing a prior construction
of the register itself requires the coordinatewise action of its selection,
transport and update operators. That dependency is kept
separate from the rigidity theorem, the finite emission catalogue and
the subsequent analytic realizations.

\subsection{Reproduction certificate and location of the data}

The program \texttt{verificar\_nucleo\_finito.py} traverses the $324$
conditions of each cursor and forms the product of the resulting classes.
The documentary catalogue is loaded after this construction to
compare emissions and their multiplicities. In the second part,
the four transition matrices are extracted from the verbatim
copy of the source; solutions are computed by exact elimination
modulo $3$ before comparison with the printed matrices. The receipt
\texttt{NUCLEO\_FINITO.json} preserves the domains, preimages,
emissions, $16$ calendars and $144$ modifications.

Verification uses exact integers and residues. It receives no
metrological values or real expansions of the constants. Its scope
is that of the finite operations and starting data specified
in this appendix. The proofs of compatibility at unbounded
depth and physical identifications retain their respective
domains and expository obligations.
